\def\stat{zatsman}



\def\tit{НАУЧНАЯ ПАРАДИГМА ИНФОРМАТИКИ: КЛАССИФИКАЦИЯ ТРАНСФОРМАЦИЙ 

ОБЪЕКТОВ ПРЕДМЕТНОЙ ОБЛАСТИ$^*$}



\def\titkol{Научная парадигма информатики: классификация трансформаций 

объектов} % предметной области}



\def\aut{И.\,М.~Зацман$^1$}



\def\autkol{И.\,М.~Зацман}



\titel{\tit}{\aut}{\autkol}{\titkol}



\index{Зацман И.\,М.}

\index{Zatsman I.\,M.}



{\renewcommand{\thefootnote}{\fnsymbol{footnote}}

\footnotetext[1]

{Исследование выполнено с~использованием инфраструктуры Цент\-ра коллективного пользования 

<<Высокопроизводительные вы\-чис\-ле\-ния и~большие данные>> (ЦКП <<Информатика>>) ФИЦ ИУ РАН 

(г.~Москва).}} 



\renewcommand{\thefootnote}{\arabic{footnote}}

\footnotetext[1]{Федеральный исследовательский центр <<Информатика и~управление>> Российской академии наук, 

\mbox{izatsman@yandex.ru}}





 \label{st\stat}



 \thispagestyle{headings}

 

\vspace*{-12pt}









  \Abst{Дано описание первых результатов создания классификации трансформаций 

объектов предметной области информатики как со\-став\-ля\-ющей системы научного знания, 

охватывающей широкий спектр информационных и~компьютерных наук. Границы ее 

предметной области, которая таким образом значительно расширена, определяются в~рамках 

концепции полиадического компьютинга Пола Розенблума. Все сущности информатики 

в~статье разделены на два глобальных класса: объекты и~их трансформации. Для каждого 

такого класса в~процессе создания научной парадигмы информатики конструируется своя 

классификация, с~формирования этих классификаций и~началось ее по\-стро\-ение. В~данной 

статье рассматриваются два верхних уровня только классификации трансформаций 

объектов предметной области. Основанием для построения самого верхнего уровня служит 

предлагаемое деление предметной области информатики на среды (ментальная, сенсорно 

воспринимаемая, цифровая и~ряд других сред), каждая из которых включает объекты одной 

природы. При этом известны примеры появления в~предметной области информатики ранее 

не использовавшихся сред. Так, с~созданием информационных технологий (ИТ), применяющих 

интерфейсы <<мозг--компь\-ю\-тер>>, в~предметную область логичным было бы включение 

объектов нейросреды, а с~разработкой методов и~средств для длительного хранения данных 

большого объема с~использованием синтезированных цепочек ДНК~---  

объектов ДНК-сре\-ды. Основанием для построения следующего уровня классификации 

трансформаций объектов служит типология знаковых систем А.~Со\-ло\-мо\-ни\-ка. Цель \mbox{статьи} 

состоит в~описании, с~одной стороны, двух верх\-них уровней классификации трансформаций 

объектов информатики, с~другой стороны, метода генерации таб\-лиц, де\-та\-ли\-зи\-ру\-ющих эту 

классификацию, что дает возможность соотнести категории трансформаций объектов 

в~таб\-ли\-цах с~интерфейсами, реализованными в~сис\-те\-мах и~средствах информатики. 

Сгенерированные таблицы включают категории трансформаций, которым соответствуют 

как уже известные интерфейсы, так и~тео\-ре\-ти\-че\-ски воз\-мож\-ные, но еще не реализованные 

в~сис\-те\-мах и~средствах информатики.}

  

\KW{объекты предметной области; трансформации объектов; классификация; среда 

предметной области; парадигма информатики}



\DOI{10.14357\!/\!08696527230412}{ZIKUWO}  



\vskip 10pt plus 9pt minus 6pt



 \thispagestyle{headings}

 

 %\vspace*{-16pt}



\section{Введение}



%\vspace*{-9pt}



  На первых этапах становления информатики объектами ее исследований 

были в~основном компьютеры и~окружающие их явления. Например, 

в~1967~г.\ А.~Ньюэлл, А.\,Дж.~Перлис и~Х.\,А.~Саймон писали: <<Везде, где 

есть феномены, может существовать наука для их описания 

и~объяснения$\ldots$ Феномены по\-рож\-да\-ют науки. Есть компьютеры. 

Следовательно, информатика [в оригинале~--- computer science] предназначена 

для описания и~изучения компьютеров. Явления, окру\-жа\-ющие компьютеры, 

разнообразны, слож\-ны и~ценны>>~[1]. Подход к~определению информатики 

как компьютерной науки без учета влияния человека на результаты работы 

сис\-тем и~средств информатики долгие годы был до\-ми\-ни\-ру\-ющим, в~том чис\-ле 

в~сис\-те\-ме компьютерного образования. Например, в~вузовском учебнике по 

информатике мож\-но было встретить такое утверж\-де\-ние: <<Для информатики 

как технической науки понятие информации не может основываться на таких 

антропоцентрических понятиях, как знание>>~\cite[c.~17]{2-zac}. В~контексте 

позиционирования информатики как технической науки в~ее предметной среде 

неявно учитывались объекты только двух сред: сенсорно воспринимаемой 

и~циф\-ро\-вой, охва\-ты\-ва\-ющей компьютерные коды, а~ментальные объекты не 

принимались во внимание.

  {\looseness=-1

  

  }

  

  В 1986~г.\ Кристен Нюгор (один из создателей языка программирования 

Симула) предложил концептуально иной поход к~определению 

информатики~[3]. Второй раздел <<The definition of the term ``informatics''>> 

его доклада на конгрессе IFIP был посвящен определению термина 

<<информатика>>. Сначала в~этом разделе Кристен Нюгор говорит:

    <<Термин ``computer science'' следует заменить на ``informatics''. Несколько 

лет назад выбор между этими двумя терминами казался, скорее всего, 

несущественным. Об\-суж\-де\-ния терминологии час\-то считают праздными, но 

иногда они могут отражать ключевые различия во мнениях или по крайней 

мере существенные акценты. Сегодня, к~сожалению, используется термин 

``computer science''. Этот термин имеет тенденцию поддерживать слиш\-ком 

узкое пред\-став\-ле\-ние об информационных сис\-те\-мах, которые в~настоящее 

время интегрируют коллективы людей и~разнообразные средства обработки 

информации, взаимодействующие как посредством межличностных связей, так и~посредством (все воз\-рас\-та\-ющей доли) электронных каналов связи>>~[3].

  

  Затем, после констатации важности четкого определения терминов 

<<информатика>> и~<<информация>>, Кристен Нюгор предлагает свою 

дефиницию первого термина:

    <<Информатика~--- это наука, которая имеет своей областью [исследований] 

информационные процессы и~\textit{связанные с~ними феномены 

в~артефактах, обществе и~природе} (курсив мой~--- И.\,З.)>>~[3]. Далее 

Кристен Нюгор цитирует определение понятия <<феномен>> по словарю 

Webster 1960~г.\ (<<любой факт, обстоятельство или событие, которые 

сенсорно воспринимаются и~могут быть научно описаны или оценены>>~[4]).

  

  Предложение учитывать влияние человека имплицитно вводит 

в~предметную об\-ласть информатики объекты среды новой природы~--- 

\textit{ментальной} (например, результаты процессов генерации научного 

знания, происходящих в~сознании исследователей). Отметим, что 

\textit{сенсорно воспринимаемая среда} содержит, в~част\-ности, знаковые 

формы представления результатов процессов генерации знания после его 

деления на концепты. В~результате, согласно Кристену Нюгору, по сравнению 

с~техническим подходом к~определению информатики, число сред ее 

предметной области расширяется до трех: \textit{ментальная}, \textit{сенсорно 

воспринимая} и~\textit{циф\-ро\-вая}.

{ %\looseness=-1



}

  

  Однако в~настоящее время она не ограничена только этими тремя средами. 

С~появлением интерфейса <<мозг--компью\-тер>> были добавлены объекты 

\textit{нейросреды}. Разработка методов записи и~хранения данных 

с~использованием синтезированных цепочек ДНК обусловила добавление 

\textit{ДНК-сре\-ды} и~ее объектов. Не исключено и~дальнейшее увеличение 

числа сред разной природы в~предметной об\-ласти информатики. Это может 

произойти, например, в~том случае, когда при проектировании будущих 

ИТ встретятся сущности, которые по своей 

природе не относятся ни к~одной из сред, ранее включенных в~предметную 

об\-ласть информатики~[5]. Подобное увеличение числа сред разной природы не 

наблюдается ни в~естественных, ни в~гуманитарных науках, что существенно 

услож\-ня\-ет позиционирование информатики в~традиционной дихотомии 

сис\-те\-мы научного знания: точные и~естественные~---  

со\-ци\-аль\-но-гу\-ма\-ни\-тар\-ные науки.

{\looseness=-1



}

  

  В статье границы предметной области информатики определяются в~рамках 

концепции полиадического компьютинга Пола Розенблума~[6]. На основе 

этой концепции предпринята попытка создать научную парадигму 

информатики~[7--10] с~учетом расширения ее предметной об\-ласти, 

предложенного Розенблумом. В~процессе создания парадигмы все сущности 

ее предметной области, значительно им расширенной, были разделены на два 

глобальных класса: объектов сред предметной области и~их трансформаций. 

Для каждого из них было начато формирование своей классификации. 

В~данной статье дано описание двух верхних уровней только классификации 

трансформаций объектов предметной области информатики. Основанием для 

построения самого верхнего уровня служит деление ее предметной области на 

среды (ментальную, сенсорно воспринимаемую, цифровую и~ряд других сред), 

каждая из которых включает объекты одной природы~\cite{7-zac, 10-zac}. 

Основанием для построения сле\-ду\-юще\-го уровня классификации 

трансформации объектов служит типология знаковых систем~[11], перечень 

видов которых приводится в~следующем разделе. Цель статьи состоит 

в~описании, с~одной стороны, двух уровней классификации трансформаций 

объектов информатики, с~другой стороны, метода генерации таб\-лиц, 

детализирующих классификацию трансформаций объектов, что дает 

возможность соотнести теоретические трансформации с~практическими 

интерфейсами, реализованными в~системах и~средствах информатики.

 

  

\vspace*{-7.7pt}

  

\section{Классификации объектов и~их~трансформаций}



\vspace*{-4pt}



  В работе~\cite{7-zac} дано описание \textit{основания классификации 

объектов} предметной области информатики для случая следующих пяти ее 

сред, включающих объекты одной и~той же природы:\\[-13pt]

  \begin{enumerate}[(1)]

  \item ментальная среда~--- это совокупность когнитивных феноменов, 

используемых и\!/\!или формируемых в~процессе познания, происходящего 

в~сознании людей (феномены, формируемые как смысловые элементы знания 

в~процессах познания с~использованием знаковых систем, будем называть 

концептами);\\[-13pt]

  \item сенсорно воспринимаемая среда, которую для краткости иногда будем 

называть \textit{информационной} средой~--- это совокупность сенсорно 

воспринимаемых феноменов, находящихся вне сознания, но 

взаимодействующих с~когнитивными феноменами ментальной среды;\\[-13pt]

  \item цифровая среда~--- это совокупность компьютерных кодов;\\[-13pt]

  \item нейросреда~--- это электрические потенциалы и~магнитные поля, 

ге\-не\-ри\-ру\-емые мозгом, которые используются, например, в~ИТ управ\-ле\-ния 

ро\-бо\-ти\-зи\-ро\-ван\-ной рукой~[12] и~в других ИТ, применяющих интерфейсы  

<<мозг--компью\-тер>>;\\[-13pt]

  \item ДНК-среда~--- это совокупность цепочек РНК и~ДНК\footnote{Например, 

модели трансляции естественных ДНК, созданные микробиологами, используются в~информатике 

при разработке методов и~средств записи для длительного хранения данных большого объема 

с~использованием синтезированных цепочек ДНК.}.\\[-13pt]

  \end{enumerate}

  

  Как было отмечено выше, дальнейшее увеличение числа сред в~предметной 

области информатики может произойти, например, при проектировании ИТ 

будущих поколений. Основанием для построения верхнего уровня 

классификации трансформаций объектов служит \textit{степень разнообразия 

природы объектов}, вовлеченных в~трансформации:\\[-13pt]

  \begin{itemize}

\item  первый класс верхнего уровня классификации включает 

трансформации объектов в~пределах среды только одной природы (далее 

природные трансформации первого порядка); например, 

конкатенация трехзначных чисел~$N$ с~$2N$ и~$3N$, результат которой 

содержит~9~цифр кроме~0, может служить примером трансформации 

первого порядка~[13];\\[-13pt]

\item  второй класс включает трансформации объектов, относящихся 

к~двум средам разной природы (далее~--- природные трансформации 

второго порядка); например, компьютерное кодирование текстов 

с~помощью таб\-лиц Unicode или представление ментальных понятий 

(концептов) словами естественного языка;\\[-13pt]

\item третий и~последующие классы включают трансформации объектов, 

относящихся к~трем и~более средам разной природы (далее~--- природные 

трансформации третьего и~более высоких порядков); например, 

присвоение одного компьютерного кода (цифровая среда) понятию 

($=$\;концепту ментальной среды) и~слову как сочетанию букв 

(информационная среда), которое его выражает, что относится к~третьему 

порядку трансформаций объектов~[14].\\[-13pt]

  \end{itemize}

  

  \begin{figure} %fig1

\vspace*{1pt}

  \begin{center}

 \mbox{%

 \epsfxsize=115.319mm 

\epsfbox{zac-1.eps}

 }



\vspace*{6pt}



{\small Два верхних уровня классификации трансформаций объектов}

\end{center}

\vspace*{-9pt}

\end{figure}

  

  Основанием для построения второго уровня классификации трансформаций 

объектов служит следующая типология знаковых сис\-тем  

А.~Соломоника~\cite[c.~131]{11-zac}: естественные знаковые сис\-те\-мы, 

образные, ес\-тест\-вен\-но-язы\-ко\-в\'{ы}е, вер\-баль\-но-не\-сло\-вес\-ные 

системы записи\footnote[1]{Под системой записи понимается знаковая система, сочетающая 

вербальные знаки с~несловесными (языки нотной записи, карт, таблиц и~др.).} 

и~формализованные знаковые сис\-те\-мы\footnote[2]{В двух последующих монографиях 

А.~Соломоник делит формализованные системы на знаковые формализованные системы первого 

и~второго порядка~[15, с.~76; 16, с.~64].}, включая математические. Введем понятие 

обобщенного текста~--- это текст, который может быть создан в~любой из 

перечисленных знаковых сис\-тем. Тогда обобщенные тексты могут быть 

естественными, образными, ес\-тест\-вен\-но-язы\-ко\-в\'{ы}\-ми,  

вер\-баль\-но-не\-сло\-вес\-ны\-ми и~формализованными. Второй уровень 

классификации трансформаций в~этой статье охватывает не все виды объектов 

предметной об\-ласти информатики, а~только перечисленные~5~видов 

текстов, \textit{вовлеченных в~процесс трансформаций}.

  %

Этот уровень содержит следующие классы (см.\ рисунок):

\begin{itemize}

\item  первый класс второго уровня классификации включает 

трансформации текс\-тов только одного вида знаковой системы (далее~--- 

знаковые трансформации первого порядка), например перевод  

ес\-тест\-вен\-но-язы\-ко\-в\'{ы}х текстов с~одного языка на 

другой\footnote[3]{Отметим, что пример перевода текста с~одного языка на другой как 

знаковая трансформация \textit{первого порядка} на втором уровне классификации относится к~классу природных трансформаций \textit{второго порядка} на верхнем уровне 

классификации трансформаций объектов, что выделено двойной стрелкой на рисунке. По 

определению такая природная трансформация относится к~двум средам разной природы: 

сенсорно воспринимаемая и~цифровая среды для машинного перевода или ментальная 

и~сенсорно воспринимаемая среды для че\-ло\-ве\-ка-пе\-ре\-вод\-чика.};

\item  второй класс включает трансформации текстов, относящихся 

к~двум видам знаковых сис\-тем (далее~--- знаковые трансформации 

второго порядка), например формирование в~редакторе Word на основе 

параллельных текстов на двух естественных языках (естественно-языков\'{ы}е знаковые системы) таб\-ли\-цы, каждая строка которой 

содержит оригинальное предложение в~первом столбце и~его перевод во 

втором (вер\-баль\-но-не\-сло\-вес\-ная сис\-те\-ма записи);\\[-13pt]

\item третий и~последующие классы включают трансформации текстов, 

относящихся к~трем и~более видам знаковых систем (далее~--- знаковые 

трансформации третьего и~более высоких порядков); например, если 

к~формированию таб\-ли\-цы в~предыдущем примере добавить запись 

результата вычисления отношения длины оригинального предложения 

к~длине его перевода для каждой строки, то это будет представлять собой 

знаковую трансформацию\footnote[1]{Вычисление характеристик текстов будем также 

называть их знаковыми трансформациями.} третьего порядка  

(ес\-тест\-вен\-но-язы\-ко\-в\'{ы}е знаковые  

сис\-те\-мы\;$\to$\;вер\-баль\-но-не\-сло\-вес\-ная сис\-те\-ма  

за\-пи\-си\;$\to$\;ма\-те\-ма\-ти\-че\-ская запись вектора отношений длин).\\[-13pt]

\end{itemize}







  На сегодняшний день процесс построения классификации трансформаций 

объектов еще далек от завершения\footnote[2]{Природные трансформации объектов 

с~указанием их сред в~будущем планируется разместить на третьем уровне классификации 

трансформаций.}. Однако итоги частичного формирования двух верхних ее 

уровней уже применяются для решения актуальных проблем информатики:\\[-13pt]

  \begin{itemize}

  \item определение отношений между компонентами модели DIKW  

(data--information--knowledge--wisdom~---  

дан\-ные--ин\-фор\-ма\-ция--зна\-ние--муд\-рость) \cite{8-zac, 9-zac};%\\[-13pt]

  \item проектирование ИТ извлечения нового знания из данных и\!/\!или 

текстовой информации~\cite{17-zac};\\[-13pt]

  \item формализация процессов обратимой статистической обработки 

текстов~\cite{18-zac}.\\[-13pt]

  \end{itemize}

  

\vspace*{-3.5pt}

  

\section{Метод генерации таблиц}



\vspace*{-3pt}



  На основе верхнего уровня классификации трансформаций предлагается 

следующий метод генерации таблиц, детализирующих классификацию 

трансформаций. Метод предназначен для генерации таблиц категорий 

трансформаций, которым соответствуют как уже известные интерфейсы, так 

и~теоретически возможные, но еще не реализованные в~системах и~средствах 

информатики.

  

  Метод генерации опишем на примере класса <<Природные трансформации 

второго порядка>> самого верхнего уровня этой классификации. По 

определению этот класс включает трансформации объектов, относящихся 

к~двум средам разной природы. Для случая пяти сред (ментальная, сенсорно 

воспринимаемая, цифровая, нейро- и~ДНК-сре\-да) с~теоретической точки 

зрения возможны максимум 20 категорий асимметричных трансформаций (см.\ 

\textit{пять строк и~пять столбцов} таблицы, начиная со второй строки 

и~второго столбца).

  

  Если во второй строке таблицы после названия среды перечислить \mbox{четыре} 

категории трансформаций ментальных объектов в~объекты четырех остальных 

сред, в~третьей строке~--- сенсорно воспринимаемых объектов в~объекты 

четырех остальных сред, пропуская название среды, в~четвертой~--- циф\-ро\-вых 

объектов в~объекты четырех остальных сред и~т.\,д., то \mbox{получим} ячейки 

с~названиями 20~категорий асимметричных трансформаций (см.\ таб\-ли\-цу, кроме последней строки и~последнего столбца). 

Например, категория <<Ментальные объ\-ек\-ты\;$\to$\;сен\-сор\-но 

воспринимаемые объекты>> (третья ячейка второй строки) реализуется 

с~по\-мощью интерфейсов пред\-став\-ле\-ния \mbox{концептов} формами 

знаков, а~категория <<Сенсорно воспринимаемые  

объ\-ек\-ты\;$\hm\to$\;менталь\-ные объекты>> (вторая ячейка третьей строки)~--- 

с~по\-мощью интерфей\-сов интерпретации (означивания) форм знаков по 

контексту их использования.

      {%\looseness=1

      

      }

  

      

      \begin{table}\footnotesize %tabl1

\protect\hspace*{3pt}  % сдвиг таблицы по горизонтали

\centering

\hbox{ %\rotcaption{\protect\small\mdseries Категории трансформаций объектов пяти сред и~будущей новой среды 

%предметной области информатики}

\label{t2}%

%\vspace*{2ex}



\begin{sideways}

\tabcolsep=1.44pt

%\begin{tabular}{|p{88.5pt}|p{105pt}|p{74.6pt}|p{59pt}|p{57pt}|p{79.9pt}|}

\begin{tabular}{|p{22mm}| p{24mm}| p{25mm}| p{25mm}| p{24mm}| p{25mm}| p{25mm}|}

\multicolumn{7}{c}{Категории природных трансформаций второго порядка объектов пяти сред и будущей новой среды 

предметной области информатики}\\

\multicolumn{7}{c}{\ }\\[-4pt]

\hline

Категории&1. Транс\-фор\-ма\-ции в~мен\-таль\-ные объ\-ек\-ты&2. Транс\-фор\-ма\-ции в~сен\-сор\-но вос\-при\-ни\-ма\-емые 

объ\-ек\-ты&3. Транс\-фор\-ма\-ции в~циф\-ро\-вые объ\-ек\-ты&4. Трансформации в~ДНК-объ\-ек\-ты&5. Транс\-фор\-ма\-ции в~объ\-ек\-ты ней\-ро\-сре\-ды&

6. Транс\-фор\-ма\-ции в~объ\-ек\-ты но\-вой сре\-ды\\

\hline

1. Транс\-фор-\newline ма\-ции мен\-таль\-ных объ\-ек\-тов&\textbf{Мен\-таль\-ная}\newline \textbf{среда}&Мен\-таль\-ные 

объ\-ек\-ты\;$\to$\;сен\-сор\-но вос\-при\-ни\-ма\-емые объ\-ек\-ты&Мен\-таль\-ные объ\-ек\-ты\;$\to$\;циф\-ро\-вые 

объ\-ек\-ты&Мен\-таль\-ные объ\-ек\-ты\;$\to$\;ДНК-объ\-ек\-ты&Мен\-таль\-ные 

объ\-ек\-ты\;$\to$\;ней\-ро\-объ\-ек\-ты&Мен\-таль\-ные объ\-ек\-ты\;$\to$\;объ\-ек\-ты новой среды\\

\hline

2. Транс\-фор\-ма\-ции сен\-сор\-но вос\-при\-ни\-ма\-емых объ\-ек\-тов&Сен\-сор\-но вос\-при\-ни\-ма\-емые 

объ\-ек\-ты\;$\to$\;мен\-таль\-ные объ\-ек\-ты&\textbf{Сен\-сор\-но вос\-при\-ни\-ма\-емая сре\-да}&Сен\-сор\-но 

воспринимаемые объ\-ек\-ты\;$\to$\;циф\-ро\-вые объ\-ек\-ты&Сен\-сор\-но вос\-при\-ни\-ма\-емые 

объ\-ек\-ты\;$\to$\;ДНК-объ\-ек\-ты&Сен\-сор\-но вос\-при\-ни\-ма\-емые 

объ\-ек\-ты\;$\to$\;ней\-ро\-объ\-ек\-ты&Сен\-сор\-но вос\-при\-ни\-ма\-емые объ\-ек\-ты\;$\to$\;объ\-ек\-ты но\-вой 

сре\-ды\\

\hline

3. Транс\-фор\-ма\-ции циф\-ро\-вых объектов&Цифровые объ\-ек\-ты\;$\to$\;мен\-таль\-ные 

объекты&Циф\-ро\-вые объ\-ек\-ты\;$\to$\;сен\-сор\-но вос\-при\-ни\-ма\-емые объ\-ек\-ты&\textbf{Циф\-ро\-вая 

сре\-да}&Циф\-ро\-вые объ\-ек\-ты\;$\to$\;ДНК-объ\-ек\-ты&Циф\-ро\-вые 

объ\-ек\-ты\;$\to$\;ней\-ро\-объ\-ек\-ты&Циф\-ро\-вые объ\-ек\-ты\;$\to$\;объ\-ек\-ты но\-вой сре\-ды\\

\hline

4. Транс\-фор\-ма\-ции ДНК-объ\-ек\-тов&ДНК-объ\-ек\-ты\;$\to$\newline$\to$\;мен\-таль\-ные объекты&ДНК-объ\-ек\-ты\;$\to$\newline

$\to$\;сен\-сор\-но воспринимаемые объекты&

ДНК-объ\-ек\-ты\;$\to$\;циф\-ро\-вые 

объекты&\textbf{ДНК-среда}&ДНК-объ\-ек\-ты\;$\to$\;ней\-ро\-объ\-ек\-ты&ДНК-объ\-ек\-ты\;$\to$\newline

$\to$\;объ\-ек\-ты новой среды\\

\hline

5. Трансформации ней\-ро\-объ\-ек\-тов&Ней\-ро\-объ\-ек\-ты\;$\to$\newline

$\to$\;мен\-таль\-ные 

объ\-екты&Ней\-ро\-объ\-ек\-ты\;$\to$\newline

$\to$\;сен\-сор\-но воспринимаемые 

объ\-екты&Ней\-ро\-объ\-ек\-ты\;$\to$\;циф\-ро\-вые объекты&Ней\-ро\-объ\-ек\-ты\;$\to$\newline

$\to$\;ДНК-объ\-ек\-ты&\textbf{Нейросреда}&Ней\-ро\-объ\-ек\-ты\;$\to$\newline

$\to$\;объ\-ек\-ты  но\-вой среды\\

\hline

6. Трансформации объектов новой среды&Объекты новой сре\-ды\;$\to$\;мен\-таль\-ные 

объекты&Объекты но\-вой сре\-ды\;$\to$\;сен\-сор\-но вос\-при\-ни\-ма\-емые объ\-екты&Объекты но\-вой 

сре\-ды\;$\to$\;циф\-ро\-вые объ\-екты&Объекты но\-вой сре\-ды\;$\to$\;ДНК-объ\-ек\-ты&Объекты 

новой сре\-ды\;$\to$\;ней\-ро\-объ\-ек\-ты&\textbf{Будущая новая среда предметной области}\\

\hline

\end{tabular}

\end{sideways}

}

\end{table}

  

  Категория <<Сенсорно воспринимаемые объ\-ек\-ты\;$\to$\;циф\-ро\-вые 

объекты>> (чет\-вер\-тая ячейка третьей строки) реализуется, в~част\-ности, 

с~помощью интерфейсов и~таб\-лиц кодирования символов компьютерными 

кодами, а~категория <<Цифровые объ\-ек\-ты\;$\to$\;сен\-сор\-но 

воспринимаемые объекты>> (третья ячейка четвертой строки)~--- с~по\-мощью 

интерфейсов и~таблиц декодирования.

   

  Таблица содержит все 20~категорий класса <<Природные трансформации 

второго порядка>> \textit{для случая пяти сред}, но не для всех из них 

в~настоящее время реализованы интерфейсы второго порядка~[12] в~системах и~средствах 

информатики (см., например, категорию <<Цифровые об\-ъек\-ты\;$\to$\;мен\-таль\-ные 

объекты>> и~категорию <<Ментальные объ\-ек\-ты\;$\to$\;циф\-ро\-вые 

объекты>> в~таблице). 

  

  Сгенерированные таблицы содержат весь спектр категорий трансформаций второго порядка

для заданного перечня сред. Это дает возможность сопоставить эти категории 

с~интерфейсами, уже реализованными в~сис\-те\-мах и~средствах информатики 

для рассматриваемого перечня сред, и~увидеть теоретически до\-пус\-ти\-мые 

интерфейсы, которые не реализованы. При этом в~случае добавления 

в~предметную область информатики новой среды сгенерированные таб\-ли\-цы 

категорий будут отражать спектр всех теоретически возможных категорий 

трансформаций ее объектов (см.\ шестую строку и~шестой столбец таблицы).

  

  В заключение этого раздела отметим, что категориям класса <<Природные 

трансформации второго порядка>> соответствуют интерфейсы только второго 

порядка, а~категориям природных трансформаций третьего и~более высоких 

порядков~--- интерфейсы третьего и~более высоких порядков~\cite{12-zac}. 

Соответственно, размерность таблицы категорий, сгенерированной на основе 

класса <<Природные трансформации третьего порядка>>, будет равна трем, т.\,е.\ с~увеличением порядка природных трансформаций растет и~размерность 

соответствующих таблиц категорий\footnote{Далеко не всегда трансформации третьего 

и~более высоких порядков можно рассматривать как последовательность трансформаций второго 

порядка, например трансформации в~процессе обучения пациента пользованию роботизированной 

руки (личностные концепты пациента, релевантные его намерениям, сигналы активности мозга как 

объекты нейросреды и~компьютерные коды)~\cite{12-zac}.}.







\section{Заключение}



  Как отмечено выше, создание научной парадигмы информатики было начато 

с~построения классификации объектов ее предметной  

об\-ласти~\cite{7-zac, 10-zac} и~классификации их трансформаций с~частичным 

описанием двух ее верхних уровней в~этой статье. Для создания парадигмы за 

основу была взята структура научной парадигмы, которую разработал 

А.~Соломоник для любой <<зрелой>> науки. Обе эти классификации относятся 

к~одному из ее восьми компонентов, которые, согласно А.~Соломонику, могут 

разрабатываться отдельно, но объединяются в~единую и~цельную 

конструкцию~\cite{15-zac, 16-zac}, а~именно: философские основы, предмет 

изуче\-ния, методы изуче\-ния, аксиоматика, классификации науки, сис\-те\-ма 

терминов, языки [знаковые сис\-те\-мы] науки и~методы верификации результатов. 

Сам тер\-мин <<научная парадигма>> трактуется А.~Соломоником 

в~соответствии с~тео\-ри\-ей Т.~Куна, которая описывает процесс смены научных 

парадигм~\cite{19-zac}. При этом А.~Соломоник отмечает тот факт, что в~книге 

Куна нет ответа на вопрос: <<Из чего долж\-на со\-сто\-ять парадигма ``зрелой'' 

науки?>>

  

  Таблицы категорий трансформаций объектов предметной области 

информатики устанавливают зависимость спектра трансформаций от перечня 

сред, к~которым принадлежат взаимодействующие объекты. Как было 

отмечено, ячейки последней строки и~столбца в~таблице предсказывают спектр 

всех теоретически возможных категорий трансформаций объектов новой среды 

при ее добавлении в~предметную область информатики. Другими словами, 

в~случае включения новой среды сгенерированные таблицы дадут весь спектр 

категорий трансформаций объектов новой среды в~объекты других сред 

и~обратно.

  

  В заключение отметим, что создание парадигмы информатики находится на 

первоначальной стадии, на которой сейчас формируется только один из восьми 

компонентов~--- классификации информатики. Для объединения компонентов 

парадигмы в~единую и~цельную конструкцию планируется на следующей 

стадии пересмотреть систему аксиом информатики как фундаментальной 

науки~\cite{20-zac}, описав \textit{в~единой сис\-те\-ме терминов} 

классификацию объектов информатики, классификацию их трансформаций, 

свойственные ей методы, включая метод  

ин\-фор\-ма\-ци\-он\-но-ма\-те\-ма\-ти\-че\-ских трансформаций~\cite{18-zac}, 

а~также сис\-те\-му ее аксиом.

  

{\small\frenchspacing

{\baselineskip=10.5pt

\addcontentsline{toc}{section}{Литература}

\begin{thebibliography}{99}

\bibitem{1-zac}

\Au{Newell A., Perlis~A., Simon~H.} Computer science~/\!\!/ Science, 1967. Vol.~157. No.\,3795. 

P.~1373--1374.

\bibitem{2-zac}

\Au{Симонович С.\,В.} Информатика. Базовый курс.~--- СПб.: Питер, 2011. 640~с.

\bibitem{3-zac}

\Au{Nygaard K.} Program development as a~social activity~/\!\!/ Information processing~/ Ed. H.-J.~Kugler.~--- North Holland: Elsevier 

Science Publishers B.V., IFIP, 1986. P.~189--198.

\bibitem{4-zac}

Webster's new world dictionary of the American language.~--- New York, NY, USA: The World 

Publishing Co., 1960. 1724~p.

\bibitem{5-zac}

\Au{Зацман И.\,М.} Таблица интерфейсов информатики как  

ин\-фор\-ма\-ци\-он\-но-компью\-тер\-ной науки~/\!\!/ На\-уч\-но-тех\-ни\-че\-ская 

информация. Сер.~1: Организация и~методика информационной работы, 2014. №\,11. 

 С.~1--15.

\bibitem{6-zac}

\Au{Rosenbloom P.\,S.} On computing: The fourth great scientific domain.~--- Cambridge, MA, 

USA: MIT Press, 2013. 307~p.



\bibitem{10-zac} %7

\Au{Зацман И.\,М.} Теоретические основания компьютерного образования: среды предметной 

области информатики как основание классификации ее объектов~/\!\!/ Системы и~средства 

информатики, 2022. Т.~32. №\,4. С.~77--89. doi: 10.14357\!/ 08696527220408.



\bibitem{7-zac} %8

\Au{Зацман И.\,М.} О научной парадигме информатики: верхний уровень классификации 

объектов ее предметной области~/\!\!/ Информатика и~её применения, 2022. Т.~16. Вып.~4. 

С.~108--114. doi: 10.14357\!/\!19922264220411.

\bibitem{8-zac} %9

\Au{Зацман И.\,М.} Данные, информация и~знание в~научной парадигме 

информатики~/\!\!/ Информатика и~её применения, 2023. Т.~17. Вып.~1. С.~116--125. doi: 

10.14357\!/\!19922264230115.

\bibitem{9-zac} %10

\Au{Зацман И.\,М.} Трансформация иерархии Акоффа в~научной парадигме 

информатики~/\!\!/ Информатика и~её применения, 2023. Т.~17. Вып.~3. С.~107--113. doi: 

10.14357\!/\!19922264230315.



\bibitem{11-zac}

\Au{Соломоник А.\,Б.}  Философия знаковых систем и~язык.~--- М.: ЛКИ, 2011. 408~с.

\bibitem{12-zac}

\Au{Зацман И.\,М.} Интерфейсы третьего порядка в~информатике~/\!\!/ Информатика и~её 

применения, 2019. Т.~13. Вып.~3. С.~82--89. doi: 10.14357\!/\!19922264190312.

\bibitem{13-zac}

\Au{Chen L., Zaharia~M., Zou~J.} How is ChatGPT's behavior changing over time?~--- Cornell 

University, 2023. 23~p. arXiv:2307.09009 [cs.CL].

\bibitem{14-zac}

\Au{Зацман И.\,М.} Кодирование концептов в~цифровой среде~/\!\!/ Информатика и~её 

применения, 2019. Т.~13. Вып.~4. С.~97--106. doi: 10.14357\!/\!19922264190416.

\bibitem{15-zac}

\Au{Соломоник А.\,Б.} Опыт современной теории познания.~--- СПб.: Алетейя, 2019. 232~с.

\bibitem{16-zac}

\Au{Solomonick A.\,B.} The modern theory of cognition.~--- Newcastle, U.K.: Cambridge 

Scholars Publishing, 2021. 242~p.

\bibitem{17-zac}

\Au{Zatsman I}. Digital spiral model of knowledge creation and encoding its dynamics~/\!\!/ 18th 

Forum (International) on Knowledge Asset Dynamics Proceedings.~--- Matera, Italy: Arts for 

Business Institute, 2023. P.~581--596. {\sf 

https://www.\linebreak researchgate.net/publication/371303696\_Digital\_Spiral\_Model\_of\_Knowledge\_\linebreak Creation\_and\_Encoding\_its\_Dynamics}.

\bibitem{18-zac}

\Au{Вакуленко В.\,В., Зацман И.\,М.} Формализованное описание статистической обработки 

информации в~базах данных~/\!\!/ Информатика и~её применения, 2023. Т.~17. Вып.~3.  

С.~93--99. doi: 10.14357\!/\!19922264230313.

\bibitem{19-zac}

\Au{Кун Т.} Структура научных революций~/ Пер. c~англ.~--- М.: Прогресс, 1977. 302~с. 

(\Au{Kuhn~T.} The structure of scientific revolutions.~--- Chicago, IL, USA: University of 

Chicago Press, 1962. 264~p.).

\bibitem{20-zac}

\Au{Зацман И.\,М.} Система семиотических аксиом информатики как фундаментальной 

науки~/\!\!/ Проб\-ле\-мы и~методы информатики: II~Научная сессия ИПИ РАН: Тезисы 

докладов.~--- М: ИПИ РАН, 2005. С.~42--44.

   

   

      \end{thebibliography}

} }







\vspace*{-6pt}



\hfill{\small\textit{Поступила в~редакцию 10.09.23}}



%\vspace*{8pt}



%\hrule



%\vspace*{2pt}







%\hrule



%\vspace*{8pt}



\newpage



\vspace*{-28pt}



\def\tit{SCIENTIFIC PARADIGM OF~INFORMATICS: TRANSFORMATION~CLASSIFICATION~OF~DOMAIN OBJECTS}







\def\titkol{Scientific paradigm of~informatics: Transformation classification 

of~domain objects}



\def\aut{I.\,M.~Zatsman}



\def\autkol{I.\,M.~Zatsman}



\titel{\tit}{\aut}{\autkol}{\titkol}



\vspace*{-8pt}



\noindent

Federal Research Center ``Computer Science and Control'' of the Russian Academy

of Sciences,   44-2~Vavilov Str., Moscow 119133, Russian Federation





\def\leftfootline{\small{\textbf{\thepage}

\hfill Sistemy i Sredstva Informatiki~---

Systems and Means of Informatics 2023 vol~33 no\,4}

}%

 \def\rightfootline{\small{Sistemy i Sredstva Informatiki~---

 Systems and Means of Informatics 2023 vol~33 no\,4

\hfill \textbf{\thepage}}}



\vspace*{3pt}  

     



\Abste{A description is given of the first outcomes of creating a classification of 

transformations of objects of an informatics subject domain as a~component 

of the scientific knowledge system covering a~wide range of information and computer sciences. 

The boundaries of its subject domain which is thus significantly expanded are defined within the 

framework of Paul Rosenbloom's concept of polyadic computing. All informatics entities in 

the paper are divided into two global classes: objects and their transformations. For each such class, 

in the process of creating the scientific paradigm of informatics, its own classification is 

constructed and with their formation its creation began. This paper discusses the two top levels of 

only the classification of transformations of objects of the informatics subject domain. The basis for 

constructing the highest level is the proposed division of the informatics subject domain into media 

(mental, sensory, digital, and a~number of other media), each of which includes objects of 

the same nature. At the same time, there are examples of the appearance of previously unused 

media in the informatics subject domain. Thus, with the creation of information technologies that use 

brain-computer interfaces, it would be logical to add the neuromedium to the informatics subject 

domain, and with the development of methods and means for long-term storage of large-volume 

data using synthesized DNA chains add the DNA-medium. The basis for constructing the next level 

of classification of object transformations is the Solomonick's typology of sign systems. The 

purpose of the paper is to describe, on the one hand, the two top levels of classification of 

transformations of objects of the informatics subject domain, and on the other hand, the method for 

generating tables detailing this classification which makes it possible to correlate the categories of 

object transformations in the tables with the interfaces implemented in computer systems and means. The 

generated tables include categories of transformations that correspond to both already known 

interfaces and theoretically possible ones but not yet implemented in computer systems and 

means.}



\KWE{domain objects; transformation of objects; classification; subject domain medium; 

informatics paradigm}







\DOI{10.14357\!/\!08696527230412}{ZIKUWO}  



\vspace*{-10pt}





\Ack



\vspace*{-3pt}



\noindent

The research was carried out using the infrastructure of the Shared Research Facilities ``High 

Performance Computing and Big Data'' (CKP ``Informatics'') of FRC CSC RAS (Moscow).



%\vspace*{-3pt}



\renewcommand{\bibname}{\protect\rmfamily References}

%\renewcommand{\bibname}{\large\protect\rm References}



{\small\frenchspacing

{ %\baselineskip=12pt

\addcontentsline{toc}{section}{References}

\begin{thebibliography}{99}



\bibitem{1-zac-1}

\Aue{Newell, A., A.~Perlis, and H.~Simon.} 1967. Computer science. \textit{Science} 

157(3795):1373--1374.

\bibitem{2-zac-1}

\Aue{Simonovich, S.\,V.} 2011. \textit{Informatika. Bazovyy kurs} [Computer science. Basic 

course]. Saint Petersburg: Piter Publishing House. 640~p.

\bibitem{3-zac-1}

\Aue{Nygaard, K.} 1986. Program development as a~social activity. \textit{Information 

processing}. Ed. H.-J.~Kugler. North Holland: 

Elsevier Science Publishers B.V., IFIP. 189--198.

\bibitem{4-zac-1}

Webster's new world dictionary of the American language. 1960. New York, NY: The World Publishing 

Co. 1724~p.

\bibitem{5-zac-1}

\Aue{Zatsman, I.} 2014. A table of interfaces of informatics as computer and information science. 

\textit{Scientific Technical Information Processing} 41(4):233--246.

\bibitem{6-zac-1}

\Aue{Rosenbloom, P.\,S.} 2013. \textit{On computing: The fourth great scientific domain}. 

Cambridge, MA: MIT Press. 307~p.

\bibitem{10-zac-1} %7

\Aue{Zatsman, I.\,M.} 2022. Teo\-re\-ti\-che\-skie osno\-va\-niya komp'\-yuter\-no\-go  

ob\-ra\-zo\-va\-niya: sre\-dy pred\-met\-noy ob\-lasti in\-for\-ma\-ti\-ki kak osno\-va\-nie  

klas\-si\-fi\-ka\-tsii ee ob''ek\-tov [Theoretical foundations of digital education: Subject domain 

media of informatics as the base of its objects' classification]. \textit{Sistemy i~Sredstva 

Informatiki~--- Systems and Means of Informatics} 32(4):77--89. doi: 

10.14357\!/\!08696527220408.



\bibitem{7-zac-1} %8

\Aue{Zatsman, I.} 2022. O~nauch\-noy pa\-ra\-dig\-me in\-for\-ma\-ti\-ki: verkh\-niy uro\-ven'  

klas\-si\-fi\-ka\-tsii ob''ek\-tov ee pred\-met\-noy ob\-lasti [On the scientific paradigm of 

informatics: The classification high level of its objects]. \textit{Informatika i~ee Primeneniya~--- 

Inform. Appl.} 16(4):73--79. doi: 10.14357\!/\!19922264220411.  

\bibitem{8-zac-1} %9

\Aue{Zatsman, I.\,M.} 2023. Dan\-nye, in\-for\-ma\-tsiya i~zna\-nie v~na\-uch\-noy  

pa\-ra\-dig\-me in\-for\-ma\-ti\-ki [On the scientific paradigm of informatics: Data, information, and 

knowledge]. \textit{Informatika i~ee Primeneniya~--- Inform. Appl.} 17(1):116--125. doi: 

10.14357\!/ 19922264230115.

\bibitem{9-zac-1} %10

\Aue{Zatsman, I.\,M.} 2023. Trans\-for\-ma\-tsiya ierar\-khii Akof\-fa v~na\-uch\-noy  

pa\-ra\-dig\-me in\-for\-ma\-ti\-ki [Transformation of the Ackoff's hierarchy in the scientific 

paradigm of informatics]. \textit{Informatika i~ee Primeneniya~--- Inform. Appl.} 17(3):107--113. 

doi: 10.14357\!/ 19922264230315.



\bibitem{11-zac-1}

\Aue{Solomonick, A.\,B.} 2011. \textit{Fi\-lo\-so\-fiya zna\-ko\-vykh sis\-tem i~yazyk} [Philosophy 

of sign systems and language]. Moscow: LKI. 408~p.

\bibitem{12-zac-1}

\Aue{Zatsman, I\, M}. 2019. Interfeysy tret'e\-go po\-ryad\-ka v~in\-for\-ma\-ti\-ke [Third-order 

interfaces in informatics]. \textit{Informatika i~ee Primeneniya~--- Inform. Appl}. 13(3):82--89. 

doi: 10.14357\!/\!19922264190312.

\bibitem{13-zac-1}

\Aue{Chen, L., M.~Zaharia, and J.~Zou.} 2023. How is ChatGPT's behavior changing over time? 

\textit{arXiv.org}. 23~p. Available at: {\sf https://arxiv.org/abs/2307.09009} (accessed October~23, 

2023).

\bibitem{14-zac-1}

\Aue{Zatsman, I.\,M.} 2019. Kodirovanie kontseptov v~tsif\-ro\-voy sre\-de [Digital encoding of 

concepts]. \textit{Informatika i~ee Primeneniya~--- Inform. Appl.} 13(4):97--106. doi: 

10.14357\!/\!19922264190416.

\bibitem{15-zac-1}

\Aue{Solomonick, A.\,B.} 2019. \textit{Opyt sovremennoy teo\-rii poz\-na\-niya} [Experience of 

modern theory of cognition]. Saint Petersburg: Aletheia. 232~p.

\bibitem{16-zac-1}

\Aue{Solomonick, A.\,B.} 2021. \textit{The modern theory of cognition}. Newcastle, U.K.: 

Cambridge Scholars Publishing. 242~p.

\bibitem{17-zac-1}

\Aue{Zatsman, I.} 2023. Digital spiral model of knowledge creation and encoding its dynamics. 

\textit{18th Forum (International) on Knowledge Asset Dynamics Proceedings}. Matera, Italy: Arts 

for Business Institute. 581--596.

Available at:

{\sf 

https://www.\linebreak researchgate.net/publication/371303696\_Digital\_Spiral\_Model\_of\_Knowledge\_\linebreak Creation\_and\_Encoding\_its\_Dynamics}

(accessed November~16, 2023).

\bibitem{18-zac-1}

\Aue{Vakulenko, V.\,V., and I.\,M.~Zatsman.} 2023. For\-ma\-li\-zo\-van\-noe opi\-sa\-nie  

sta\-ti\-sti\-che\-skoy ob\-ra\-bot\-ki in\-for\-ma\-tsii v~ba\-zakh dan\-nykh [Formalized description 

of statistical information processing in databases]. \textit{Informatika i~ee Primeneniya~--- Inform. 

Appl.} 17(3):93--99. doi: 10.14357/19922264230313.

\bibitem{19-zac-1}

\Aue{Kuhn, T.} 1962. \textit{The structure of scientific revolutions}. Chicago, IL: University of 

Chicago Press. 264~p.

\bibitem{20-zac-1}

\Aue{Zatsman, I.\,M.} 2005. Sistema se\-mio\-ti\-che\-skikh ak\-si\-om in\-for\-ma\-ti\-ki kak  

fun\-da\-men\-tal'\-noy nau\-ki [System of semiotic axioms of informatics as a~fundamental 

science]. \textit{Problemy i~metody informatiki: II Nauchnaya sessiya IPI RAN. Tezisy dokladov} 

[Problems and Methods of Informatics: 2nd Scientific Session of the IPI RAS. Abstracts]. Moscow: IPI RAN. 

42--44.



\end{thebibliography}

} }



\vspace*{-6pt}



\hfill{\small\textit{Received September 10, 2023}} 



\vspace*{-14pt} 



\Contrl



\vspace*{-3pt}

\noindent

\textbf{Zatsman Igor M.} (b.\ 1952)~--- Doctor of Science in technology, head of department, 

Federal Research Center ``Computer Science and Control'' of the Russian Academy of Sciences,  

44-2~Vavilov Str., Moscow 119333, Russian Federation; \mbox{izatsman@yandex.ru}

  

\label{end\stat}



\renewcommand{\bibname}{\protect\rm Литература}    

